%!TeX program = pdflatex
% =====================================================================
%  HercUNet — combined submission paper (Stage 1 + Stage 2)
%  Self-contained: compiles with a stock pdflatex / Overleaf, no special
%  document class required. Figures are read from ../images/ (submission/images).
%  To match Socratic's ACM-TOG look on Overleaf, swap the \documentclass line
%  for  \documentclass[acmtog,review,anonymous]{acmart}  and adapt the front
%  matter; everything below is class-agnostic.
% =====================================================================
\documentclass[10pt,twocolumn,a4paper]{article}

\usepackage[a4paper,top=0.7in,bottom=0.7in,left=0.62in,right=0.62in,columnsep=16pt]{geometry}
\usepackage{microtype}                                  % tighter, cleaner line-packing
\usepackage{amsmath}
\usepackage{amssymb}
\usepackage{mathtools}
\usepackage{upgreek}
\usepackage{graphicx}
\usepackage{xcolor}
\usepackage{booktabs}
\usepackage[font=footnotesize,labelfont=bf]{caption}
\usepackage{enumitem}
\usepackage{stfloats}                                   % allow figure* at column bottoms in twocolumn
\usepackage[colorlinks=true,linkcolor=black,citecolor=blue,urlcolor=blue]{hyperref}

% loosen float placement so full-width figures distribute instead of piling up
\renewcommand{\topfraction}{0.92}
\renewcommand{\bottomfraction}{0.7}
\renewcommand{\textfraction}{0.06}
\renewcommand{\floatpagefraction}{0.75}
\renewcommand{\dbltopfraction}{0.92}
\renewcommand{\dblfloatpagefraction}{0.75}
\setcounter{topnumber}{3}
\setcounter{bottomnumber}{2}
\setcounter{totalnumber}{5}
\setcounter{dbltopnumber}{3}
% tighten vertical rhythm a touch
\setlength{\parskip}{0pt}
\setlength{\abovedisplayskip}{4pt}
\setlength{\belowdisplayskip}{4pt}

\graphicspath{{../images/}{images/}{./}}

% --- convenience macros --------------------------------------------------
\newcommand{\um}{\ensuremath{\,\upmu\mathrm{m}}}      % micrometres
\newcommand{\nhat}{\ensuremath{\hat{\mathbf{n}}}}
\newcommand{\code}[1]{\texttt{#1}}
\newcommand{\dgrad}{\ensuremath{\nabla\phi}}

% figure-with-missing-asset fallback (so the doc builds before all assets exist)
\newcommand{\figasset}[3]{% {file}{width}{placeholder-height}
  \IfFileExists{../images/#1}{\includegraphics[width=#2]{#1}}%
    {\fbox{\parbox[c][#3][c]{#2}{\centering\small\textbf{FIGURE RESERVED}\\[.4em]\texttt{\detokenize{#1}}}}}}

\begin{document}

\twocolumn[{%
\centering
{\LARGE\bfseries HercUNet\par}\vspace{4pt}
{\large A self-refining UNet architecture for sheet detection in Herculaneum
scrolls\par}\vspace{9pt}
{\large Jaime Lomeli-R.\quad\textnormal{\normalsize\textbullet}\quad\today\par}\vspace{6pt}
}]

{\small
\begin{center}\textbf{Abstract}\end{center}
\noindent
Sheet detection must not only find surface but \emph{separate} touching wraps.
Existing detectors, trained per corpus, do not generalise: they merge touching
wraps and break thin sheets---topological errors voxel-accuracy metrics never
expose. We present a two-stage pipeline. \textbf{Stage~1} generates surface
pseudo-labels fully unsupervised, at a quality high enough to train an nnU-Net,
on the premise that separation is a problem of \emph{measurement}, not learning:
adjacent wraps are locally identical, and the only cue between them is the thin
gap seen edge-on across the sheet normal. From the raw micro-CT we estimate a
per-voxel orientation frame, grow oriented \emph{primitives} longer than the
inter-wrap spacing, measure same-sheet and across-gap relations with an
anisotropic curvature-following slab, and train a contrastive embedding in which
wraps form an ordered ladder; a size-penalised density clustering recovers
per-sheet instances, rasterised into a surface crest and a per-voxel sheet
identity with honest confidence. \textbf{Stage~2} distils these into an iterative
refiner $D(\mathrm{CT},\mathrm{prev})\!\to\!\mathrm{surface}$ applied repeatedly on
its own output: a carried CT orientation field supplies the separating cue that
survives fusion, a constrained-MALIS affinity objective decouples completing a
sheet from merging two, and the iteration is taught by a synthetic prior channel
and online DAgger. From learned tractography we take one idea---that a network can
track a structure through the raw signal beyond the traces it was shown---to grow
sheets into material the labels never reached. We are candid about the open limit
(separating \emph{already-merged} wraps is unsolved) and about evaluation: we
report no Dice, and explain why.
\par}

\medskip
\noindent\textbf{Overview.} The pipeline runs over one shared streaming data layer
for teravoxel OME-Zarr volumes, in two stages: Part~\ref{part:labels} is the
unsupervised pseudo-label generator (\emph{HercuLabels}); Part~\ref{part:refiner}
the learned refiner (\emph{HercUNet}) that turns those labels into a dense,
self-correcting, cross-volume predictor. \emph{A note before comparing:} this
detector predicts sheet \textbf{medials}, a few voxels inboard of the
\textbf{faces} every other label (manual or mined from m7) annotates---a
deliberate, different target, not an error, and the main reason we report no Dice
against face labels (\S\ref{sec:eval}). \emph{It is better to have a detection
whose face lives flat at a few voxels' offset than one flat at zero offset only in
some regions and waving harshly everywhere}: a consistent, correctable offset beats
intermittent exactness (Figure~\ref{fig:facemed}).

\begin{figure}[t]
  \centering
  \figasset{faces-vs-medials-1447-yz.jpg}{\linewidth}{4.5cm}
  \caption{Faces vs.\ medials --- a YZ cross-section of PHerc1447. m7 annotates
  the sheet \textbf{face} (blue); our detector predicts the \textbf{medial} crest
  (magenta), sitting a consistent few voxels inboard of the same sheet. The
  offset is constant and trivially correctable, not a miss.}
  \label{fig:facemed}
\end{figure}

\begin{figure*}[t]
  \centering
  \figasset{banner-PHerc1447-m7-passes.jpg}{\textwidth}{4cm}
  \caption{The iterative refiner across passes (magenta) versus the m7 baseline
  (blue), on a 2\,mm cross-section of PHerc1447.}
  \label{fig:banner}
\end{figure*}

% =====================================================================
\part{Pseudo-Label Generation (HercuLabels)}
\label{part:labels}
% =====================================================================

We generate pseudo-labels fully unsupervised, at a quality high enough to train an
nnU-Net. The central difficulty is \emph{measurement}, not learning: inside a small
window two adjacent wraps are locally indistinguishable---same fibre texture,
thickness, appearance---and the \emph{only} cue between them is the thin gap seen
edge-on across the sheet normal. Any method that decides membership from a single
scalar (a density, a coordinate) is fooled wherever sheets bend, fold, or compress
together.

Our pipeline is built around that fact. From the raw micro-CT we estimate a
per-voxel orientation \textbf{frame field}; along it we grow oriented surface
\textbf{primitives} longer than the inter-wrap spacing; between primitives we
measure two relations---\emph{same-sheet} affinity and \emph{across-gap}
opposition---with an \textbf{anisotropic, curvature-following, asymmetric slab}
that conforms to the local material. These train a low-dimensional
\textbf{contrastive embedding} in which same-sheet primitives collapse and
adjacent wraps separate into an ordered ladder; a \textbf{size-penalised density
clustering} then recovers per-sheet instances, rasterised into the two fields
the refiner consumes---a surface crest and a per-voxel \textbf{sheet-instance
identity}, each with an honest confidence. The principle throughout: clean
sheets with truthfully-marked uncertain regions beat confidently-wrong ones.

\section{Setting and notation}
We work on one window at a time, usually a cube of 192 voxels. A direction always
means the sheet \textbf{normal} \nhat---across the sheet, along which neighbouring
wraps separate---and \emph{in-plane} the two tangent directions.

\section{The cleaned substrate and the orientation frame}
Every later stage needs two things at every point: a clean local density, and
the direction the sheet faces. Both come from one operator---the
\textbf{structure tensor} of the density---used first to clean the volume and
then to read the frame off it.

\paragraph{The structure tensor.} At each voxel we differentiate the density
with a derivative-of-Gaussian at a small \emph{noise} scale (so it does not
amplify the voxel grain) and take the outer product of the gradient---a
$3\times3$ matrix. A lone outer product is rank-one and noisy; the structure
tensor \textbf{averages} them over a larger \emph{integration} neighbourhood (a
second Gaussian at the sheet scale). That averaging is decisive: individual
gradients on a fibre-textured sheet scatter, but their averaged outer product
has a dominant eigenvector pointing reliably \textbf{across} the sheet.
Diagonalising (ascending eigenvalues) gives an orthonormal frame---the
\textbf{largest}-eigenvalue eigenvector is the \textbf{normal} \nhat, the
\textbf{smallest} the in-plane \textbf{fibre} direction, and the eigenvalue
anisotropy is a \textbf{coherence} (near one on a crisp sheet, near zero in noise
or where orientations collide). The two scales---small derivative, larger
integration---are what let a single robust orientation survive the fibre texture.

\paragraph{Cleaning the substrate (coherence-enhancing diffusion).} The raw CT is
grainy and its gaps faint, degrading every geometric estimate. We pre-clean it
with Weickert-style \textbf{coherence-enhancing anisotropic diffusion}~\cite{weickert}
steered by the frame: intensity flows \textbf{freely in-plane and barely across
the normal} (diffusion tensor $=$ identity with its across-normal eigenvalue
$\ll1$), so the flux is the gradient minus its across-normal part. Over a few
explicit steps---the frame re-estimated as edges sharpen---this smears each sheet
into a clean connected band while \textbf{widening and sharpening the inter-wrap
gaps}, the two properties the slab and gap trace depend on. It is destructive by
design and used only as this geometric substrate, never as a detection input;
the frame field is the structure tensor of \emph{this cleaned density}.

\paragraph{Why the frame is read from the CT, and sign-free.} The frame is
derived, never learned, and read from the CT rather than any prediction: where
two wraps have visually merged the CT still carries \emph{both} orientations even
though a detector fused them, and that surviving discontinuity is the cue that
later separates them. The normal is \textbf{sign-free}---within a window nothing
says which way the umbilicus lies, so any quantity flipping under
$\nhat\to-\nhat$ is avoided throughout.

\section{The primitive: an oriented 2.5-D surface patch}
We do not classify voxels. We tile each sheet with a \textbf{primitive}---a small
planar cross of a \textbf{spine} and perpendicular \textbf{ribs}---sampling the
local tangent frame. Each arm is a \textbf{streamline}: an integral curve of a
direction field traced from a seed by \textbf{fixed-step midpoint (RK2)
integration}, growing both ways, sampled by trilinear interpolation. What
distinguishes ours from fibre tractography is \emph{which} field is integrated
and how it is constrained.

Order the structure-tensor eigenvectors by eigenvalue: the \textbf{normal}
$\mathbf{e}_0$ (largest, across the sheet), the \textbf{fibre} $\mathbf{e}_2$
(smallest, along the papyrus fibres), and the \textbf{middle} $\mathbf{e}_1$ (the
second in-plane tangent, across the fibres). Each arm is grown aligned to one
in-plane axis.

\paragraph{Growing the spine.} A spine is seeded on high-coherence material on a
40\um\ grid and grown \textbf{along the fibre} $\mathbf{e}_2$. The trick is to
\textbf{freeze a plane at the seed}---spanned by the seed's normal $\mathbf{e}_0$
and fibre $\mathbf{e}_2$, so its plane-normal is the middle eigenvector
$\mathbf{e}_1$---and require the spine to stay both on the sheet and in that
plane. The step direction is the local normal crossed with the frozen
plane-normal,
\begin{equation}
  \mathbf{d}(p) \;=\; \mathrm{localNormal}(p)\times \mathbf{e}_1(\text{seed}),
\end{equation}
perpendicular to the local normal (riding the sheet) and to $\mathbf{e}_1$
(staying in the plane). Where the sheet is flat this is the fibre direction;
where it undulates, $\mathbf{d}$ \textbf{pitches} across the normal to follow the
surface but can \textbf{never yaw} out of the frozen plane---projected onto the
seed's tangent plane, the spine is a straight line along the fibre. Two details:
the step is driven by the \textbf{local normal} (the structure tensor's most
stable output) rather than the fibre eigenvector, which is ill-defined where the
two in-plane eigenvalues are close and makes a streamline wander; and each 8\um\
step is \textbf{sign-coherent} and \textbf{gap-shy} (coasting through a few
low-signal steps), so the spine drives through a locally incongruent stretch
instead of stopping. Growth halts at 250\um\ per side.

\paragraph{Growing the ribs.} Ribs are the dual. They are seeded at
\textbf{arc-length samples of the spine, every 20\um}---so each begins on a recto
point---and grown along the \textbf{middle eigenvector} $\mathbf{e}_1$
(perpendicular to the spine), same 8\um\ step, out to 100\um\ per side. The rule
is identical with the frozen plane-normal swapped from $\mathbf{e}_1$ to the
\textbf{fibre} $\mathbf{e}_2$: the rib is flat in the fibre direction while
pitching across the normal. Ribs gate on \textbf{planarity} rather than
fibre-coherence, so they are \emph{less} shy than the spine and run straight
across low-congruency regions; each remembers its parent spine.

\paragraph{Why one direction is held flat.} Each arm grows freely along one
in-plane axis while the other is \textbf{frozen}---the spine flat to
$\mathbf{e}_1$, each rib flat to $\mathbf{e}_2$. That constraint makes the
primitive \emph{hug} the sheet: free to pitch across the normal to track
undulation, forbidden to yaw sideways (Figure~\ref{fig:meshlet}). Without it a
streamline meanders, loops, or slips across a faint gap onto the neighbour
wherever the in-plane orientation is ambiguous.

\begin{figure}[h]
  \centering
  \figasset{meshlet_primitive.jpg}{0.85\linewidth}{5cm}
  \caption{The 2.5-D meshlet primitive: a spine grown along the fibre and ribs
  grown perpendicular, each confined to a seed-frozen plane so it pitches across
  the normal but never yaws within the sheet.}
  \label{fig:meshlet}
\end{figure}

Points are sampled along the spine and ribs at 20\um\ spacing to form a small
cloud; each point is tagged by which arm it came from---spine-derived
(\textbf{recto}) or rib-derived (\textbf{verso})---a distinction the embedding
later relies on. A primitive is a piece of geometry only; it is never treated as
a node in a graph, and no connectivity is asserted between primitives at this
stage.

A window tiles into of order \textbf{twenty thousand spines} and a few hundred
thousand ribs, sampling to a few million points---of which, because ribs are
short and dense, only about \textbf{10\,\% are recto (spine) and 90\,\% verso
(rib)}, an imbalance corrected later. The length scales (40\um\ seed grid,
20\um\ point spacing, 250/100\um\ arms) were set by sweeping seed stride and
sampling distance for a quality-vs-runtime operating point that keeps generation
tractable while still resolving sheets cleanly---not a fundamental limit.

Two properties make the primitive the right unit: it is \textbf{oriented}, so the
slab can be posed in the sheet's own frame; and it is \textbf{longer than the
merge} it must survive---a patch spanning more than the inter-wrap gap cannot be
explained by a single wrong wrap, so the cue to separate two wraps lives within
the primitive, not only across many.

\section{The surface slab --- the central measurement}
All membership decisions reduce to one question about a pair of points:
\emph{does the candidate lie on the same sheet as the anchor?} We answer it with
a \textbf{slab}---a soft, thin, sheet-shaped region attached to the anchor,
scoring high inside and low outside. It is the core measurement of the method,
and its shape is what lets the pipeline follow curved and compressed sheets
without leaking onto neighbours (Figure~\ref{fig:slab}). The same machinery serves
both relations: a \textbf{positive} pulls same-sheet points together (this
section), a \textbf{negative} pushes the adjacent wrap away (\S\ref{sec:negatives}).

\begin{figure*}[tb]
  \centering
  \figasset{sheet_slab.jpg}{\textwidth}{5cm}
  \caption{The positive slab. Left: the slab is an anisotropic Gaussian hugging
  a curved sheet---wide in-plane $\sigma_s$, tight (asymmetric) $\sigma_n$ across
  the normal, curvature-following---so a far same-sheet point is captured while
  the adjacent wrap is cut off. Right: mutuality---a pair counts as positive only
  if both slabs contain each other, which rejects the adjacent wrap (A's slab may
  reach C, but C's slab, with its normal flipped, does not reach A).}
  \label{fig:slab}
\end{figure*}

\subsection{The slab: an anisotropic Gaussian in the surface frame}
Let the anchor (one of the points sampled at the spines) sit at a point with
unit normal \nhat, and let a candidate be reached by the offset vector
$\mathbf{v}$. In the anchor's local \textbf{Darboux frame}---the two in-plane
tangents together with the normal---we split $\mathbf{v}$ into a signed distance
\emph{across} the sheet and the residual \emph{in-plane} displacement:
\begin{itemize}[nosep]
  \item across-normal distance: $d_n = \mathbf{v}\cdot\nhat$;
  \item in-plane residual: $\mathbf{v}_\parallel = \mathbf{v}-d_n\nhat$, with
        in-plane distance $\rho=|\mathbf{v}_\parallel|$.
\end{itemize}
The slab score is an \textbf{anisotropic Gaussian} in this frame---$\exp(-d^2)$
for a \textbf{Mahalanobis distance} $d$, generous in-plane and tight across the
normal:
\begin{equation}
  S(\text{anchor}\!\to\!\text{candidate}) =
  \exp\!\Big(-\rho^2/\sigma_s^2 - d_n^2/\sigma_n^2\Big),
\end{equation}
with $\sigma_n\ll\sigma_s$. A point on the same flat sheet has $d_n\approx0$ at
any $\rho$ and scores near one; a point on the adjacent wrap has large $|d_n|$
and is cut off---an isotropic neighbourhood could not tell a far same-sheet point
from a near adjacent-wrap one. Two refinements follow: a curvature-corrected
distance (\S\ref{sec:curv}) and a two-sided, material-adaptive $\sigma_n$
(\S\ref{sec:asym}).

\subsection{Following the curvature}\label{sec:curv}
Sheets are not flat, and a na\"ive across-normal test punishes a curved sheet's
own far parts as if they left it. We correct with the local \textbf{shape
operator} (the Jacobian of the normal field): from it we read the normal
curvature $\kappa$ along $\mathbf{v}_\parallel$ and replace $d_n$ with a distance
measured from the osculating sheet rather than the flat tangent plane:
\begin{equation}
  r \;=\; \frac{d_n - \tfrac12\,\kappa\,\rho^2}{\sqrt{1+\kappa^2\rho^2}}.
\end{equation}
The $\tfrac12\kappa\rho^2$ term is how far the curved sheet has risen off its
tangent plane at $\rho$, so a point following the bend keeps $r\approx0$ while
only a genuine gap-crossing earns a large $r$. The slab thus hugs a rolling sheet
instead of slicing a chord.

\subsection{The asymmetric across-normal width}\label{sec:asym}
Putting the curvature-corrected $r$ in place of $d_n$ gives the score we actually
evaluate,
\begin{equation}
  S(\text{anchor}\to\text{candidate}) \;=\;
  \exp\!\Big(-\rho^2/\sigma_s^2 - r^2/\sigma_n^2\Big),
\end{equation}
with in-plane reach $\sigma_s$ ($\sim$180\um, the primitive's scale) and an
across-normal half-width $\sigma_n$ that is the delicate part. A fixed $\sigma_n$
is wrong: near a tight gap it bleeds onto the neighbour, and the material around
the medial is \textbf{not symmetric}---the distance to the ridge edge and to the
next wrap differs on the two sides.

We make the width \textbf{two-sided and material-adaptive}. Sampling density
outward along $\pm\nhat$, on each side we find the first point where it falls
below a fixed fraction of the anchor value---the local ridge edge, where the gap
begins. The half-width for that side is a fixed fraction of that edge distance,
clamped so it can only \textbf{tighten} below the default (near edges and tight
gaps), never widen (clear space falls back to the cap). The slab uses
$\sigma_n^{+}$ for $d_n\ge0$ and $\sigma_n^{-}$ otherwise: an \textbf{asymmetric
slab} conforming to the material envelope---further on the open side, pulled in on
the crowded side, never crossing a gap the density reports, on either face
independently.

\subsection{From the slab to positive (same-sheet) pairs}
An anchor's slab tells us which candidates \emph{it} considers same-sheet, but an
adjacent wrap across the normal can fall inside its generous in-plane reach. We
therefore require \textbf{mutuality}: a pair is positive only if each lies in the
other's slab, evaluated in its \emph{own} frame; the weight is the product of the
two scores. This structurally rejects the neighbour---one sheet's slab may reach
across the gap, but the neighbour's, with its normal pointing back, does not reach
the anchor, so the product collapses.

One safeguard is a \textbf{down-weighting, never a repulsion}: a soft attenuation
that fades a positive as its across-normal distance grows, weakening a spurious
glue across a small step. Being one-directional---it can stop a bad glue but never
push apart---keeps it \textbf{fold-safe}: a genuine fold is held together by its
in-plane path even as the attenuation fades.

\section{Opposition across the gap --- the negatives}\label{sec:negatives}
Pulling same-sheet points together does not separate wraps; something must
actively \textbf{push the adjacent wrap away}. Choosing these negative pairs is
the hardest part of the method; it went through three generations, and the
failure of the first motivates the third.

\subsection{The corridor (the failed first attempt)}
The first approach took, from each patch, \textbf{any} other displaced by more
than a wrap-thickness along the normal, on the logic that a same-sheet point
projects to near-zero across-normal distance. This holds on a \textbf{flat} sheet
but not a \textbf{curved} one: a bend lifts a sheet's own far parts off the
tangent plane (the $\tfrac12\kappa\rho^2$ rise), so they project large
across-normal distances and are swept up as false negatives; the embedding pushes
a sheet \textbf{apart from itself} and sheets fragment. The corridor had no notion
of a gap---it measured only across-normal distance, exactly the quantity
curvature corrupts.

\subsection{Gap-gating (the working method, used to generate the corpus)}
The fix: fire a negative \textbf{only where there is a demonstrable gap}, placed
on its far side. The rule operates on the \textbf{relative density profile} along
the normal, normalised by the anchor value so it is contrast-independent:
\begin{equation}
  \mathrm{rel}(t) \;=\; \frac{d(\text{anchor}+t\,\nhat)}{d(\text{anchor})},
  \quad t=0,2,4,\dots\um
\end{equation}
sampled on both normal sides.

\paragraph{The valley-then-rise gate.} A side fires only if three conditions hold
together: the material genuinely ends ($\sigma_n<\sigma_{\max}$); the profile
\textbf{falls into a valley}, $\mathrm{rel}(t)<\tau$ beyond a conservative gate
$t\ge\max(\mathrm{margin}_k\cdot\sigma_n,\,t_{\text{floor}})$ (so the anchor's own
tail is never read as a gap); and it \textbf{rises again},
$\mathrm{rel}(t')\ge\tau_{nb}$ at $t'>t$---\emph{material $\to$ gap $\to$
material}. With $\tau\approx0.3$, $\tau_{nb}\approx0.5$,
$t_{\text{floor}}\approx15$\um, $\mathrm{margin}_k\approx2.5$, this is a
\textbf{matched criterion for one inter-wrap gap}. If the profile continues---through
a fold, or compressed material with no valley---no negative is produced, so the
method is \textbf{fold-safe by construction}.

\paragraph{The gather point.} The distance of the neighbour rise, $D=t'$, is the
measured local \textbf{pitch} to the next wrap, and the negative is gathered at
the \textbf{far ridge}: we query the point cloud for the nearest patches to
\begin{equation}
  \mathbf{c} \;=\; \text{anchor}+D\,\nhat \qquad (\text{the blob centre}),
\end{equation}
not to the anchor. Those nearest-to-$\mathbf{c}$ patches \emph{are} the adjacent
wrap. (Gathering near the anchor instead---as one naturally would---returns the
anchor's own in-plane neighbourhood, which never reaches across the gap; that
mistake once produced essentially zero negative coverage.)

\paragraph{The weight.} Each gathered candidate $q$ is scored by a
\textbf{separable anisotropic kernel} in the normal-ray frame, decomposing
$q-\text{anchor}$ into an along-normal projection $\mathrm{proj}$ and a
perpendicular offset $\mathrm{perp}$:
\begin{equation}
  \mathrm{neg\_w}(q) \;=\; e^{-\mathrm{depth}/\sigma_d}\cdot
  e^{-\frac12(\mathrm{perp}/\sigma_{ip})^2}\cdot e^{-\mathrm{proj}/\lambda}.
\end{equation}
The three factors: a \textbf{depth gate} $e^{-\mathrm{depth}/\sigma_d}$ on how
deep the valley fell (deeper $\Rightarrow$ stronger push; a shallow same-sheet
valley $\Rightarrow\approx0$, the fold-safety); a \textbf{Gaussian across the
ray} $e^{-\frac12(\mathrm{perp}/\sigma_{ip})^2}$ keeping the push on-axis
($\sigma_{ip}\approx40$\um); and an \textbf{exponential decay along the ray}
$e^{-\mathrm{proj}/\lambda}$ favouring the nearest wrap ($\lambda\approx60$\um).
It is the dual of the positive slab---a \textbf{Gaussian} centred on the anchor
there, an anisotropic kernel centred on the \textbf{far ridge} here, Gaussian
on-axis and exponentially near-favouring.

\paragraph{Ordinality for free.} Because the trace stops at the \emph{first}
valley-then-rise, the negative lands on the \textbf{immediately} adjacent wrap
and never skips one---the property that later orders the embedding
(\S\ref{sec:embedding}). This gap-gated selection on the asymmetric slab frame
(Figure~\ref{fig:negative}) generated the entire corpus.

\begin{figure*}[tb]
  \centering
  \figasset{gap_gated_negative.jpg}{\textwidth}{5cm}
  \caption{The gap-gated negative. Left: from the anchor the density is traced
  along the normal; a negative fires only on a valley-then-rise (material $\to$
  gap $\to$ material), and the push is gathered at the far ridge
  $\mathbf{c}=\text{anchor}+D\nhat$ where the profile rises, weighted by a kernel
  that is Gaussian across the ray and exponentially near-favouring along it.
  Right: the relative profile $\mathrm{rel}(t)$ with the valley threshold $\tau$,
  the neighbour-rise threshold $\tau_{nb}$, the conservative gate, the valley
  depth, and the measured pitch $D$.}
  \label{fig:negative}
\end{figure*}

\subsection{The slab-field refinement (the current improvement)}
Gap-gating weakens exactly where wraps are crushed together: the raw density
between them stays grey rather than dropping to a valley, and the boundary test
(which requires the across-normal width below its cap) rejects the thick material
outright, so few anchors get a negative and residual merges survive.

The improvement makes the gap \textbf{visible} first. Instead of the raw density
we \textbf{splat every patch's asymmetric slab into a shared volume}. Because the
slabs are thin across the normal, the compressed inter-wrap space becomes a
\textbf{deep valley} in the accumulated field---the anisotropy manufactures
contrast the raw density lacked---and the valley test fires where it could not.
Two smaller changes accompany it (decoupling the boundary test from the tightness
cap, relaxing the valley threshold), together raising negative coverage in
compressed windows several-fold with no regression on clean ones. This is for
future label generation; the existing corpus predates it.

\subsection{The honest limit}
Where two wraps genuinely \textbf{touch}, with no valley at any scale, no density
cue can separate them---the winding continues and the normal flips across the
fold. This is topological, not a tuning failure, and no slab refinement addresses
it. Rather than emit a confident wrong boundary, the pipeline marks those regions
\textbf{low-confidence} (\S\ref{sec:confidence}) and leaves the fix to an
orientation/continuity signal outside this stage.

\section{The contrastive embedding}\label{sec:embedding}
The two relations train a small embedding assigning \textbf{one vector per
primitive} (a spine and its ribs share one) in \textbf{eight dimensions}---the
relations' effective rank is low and a compact space clusters far more cleanly. A
window carries of order \textbf{twenty thousand} vectors.

\paragraph{The training scheme.} The slab tables give each anchor a shortlist (up
to $\sim$48 positives, $\sim$32 negatives), but we sample rather than use them
all. Each step draws a large batch of anchors and, per anchor, a \textbf{few
pairs}---one recto and one verso positive ($\propto$ mutual slab weight) and one
gap negative ($\propto$ push weight)---for one Adam step. Over training each
anchor is visited $\sim$100 times, so every primitive sees many partners and
opponents without materialising the full pair graph. Anchors are the sample
points themselves ($\sim$10\,\% recto / $\sim$90\,\% verso).

\paragraph{The loss.} Four terms are summed each step, the first two the working
forces and the last two a faint regulariser borrowed from \textbf{VICReg}
(variance--invariance--covariance regularisation)~\cite{vicreg}:
\begin{itemize}[nosep]
  \item \textbf{Invariance (pull)}---the mean squared distance from the anchor
  to each of its drawn positives, $\|z_a-z_p\|^2$. Always on; it is the
  sheet-cohesion force.
  \item \textbf{Push}---a hinge on the distance to the drawn gap negative,
  $\mathrm{relu}(\mathrm{margin}-\|z_a-z_{\mathrm{neg}}\|)$, weighted by the
  negative's strength and applied only where a gap actually fired. It separates
  the adjacent wrap.
  \item \textbf{Variance}---a per-dimension floor,
  $\mathrm{relu}(1-\mathrm{std}_d)$ averaged over dimensions, which forces every
  embedding dimension to keep a minimum spread so the space cannot collapse to a
  point or a line.
  \item \textbf{Covariance}---the sum of the squared off-diagonal covariances
  between dimensions, which decorrelates the dimensions so the eight axes carry
  independent information rather than redundant copies of one.
\end{itemize}
The variance and covariance terms carry a \textbf{tiny weight} ($\sim$a tenth and
a hundredth of the pull): a stabiliser, not the objective. At healthy negative
coverage the push is its own anti-collapse; their only job is to stop a
\emph{starved} window---too few negatives fired---from imploding into a blob.
Given real weight they overwhelm pull and push and smear the sheets together,
exactly the failure of full VICReg (no negatives) as the whole loss.

Two choices matter disproportionately. First, because verso points vastly
outnumber recto ones, uniform sampling starves the recto links that tie a sheet
together in-plane; we \textbf{balance} them, drawing an equal-weight recto and
verso positive per anchor---the largest single quality gain in the embedding.
Second, wrap ordering is left to \textbf{emerge}: each negative lands on the
immediately adjacent wrap, so wrap $N$ is pushed off both $N{-}1$ and $N{+}1$ and
the embedding self-organises into a ladder where two-apart wraps sit at opposite
ends and cannot glue. Any residual merge is bounded to \emph{contiguous}
wraps---the only geometrically plausible place.

\section{Clustering: probabilistic excess-of-mass}
The embedding is clustered by density, but standard density clustering has a
systematic bias that is fatal here; the correction is a small, general
contribution.

\paragraph{Why the standard method over-merges.} Hierarchical density
clustering~\cite{hdbscan} builds a tree of nested clusters and selects a flat set
by \emph{excess of mass}: each cluster is scored by its total density-stability,
and a parent is kept over its children when its stability exceeds their sum. The
bias is structural---a child's stability accrued \textbf{before it separates} is
credited to the parent, inflating its score toward keeping the merge---and in our
data a single cluster swallows three or more wraps. The usual patch, a hard
max-size cutoff, is brittle and does not transfer across window scales.

\paragraph{The fix: a smooth size penalty on a scale-invariant split signal.} We
keep the tree but change the rule. Bottom-up, at each node we compare its
stability against its children's sum as a \textbf{normalised difference}
\begin{equation}
  \Delta \;=\; \frac{S_{\text{node}} - \sum S_{\text{children}}}
                    {|S_{\text{node}}| + |\sum S_{\text{children}}|}\;\in[-1,1].
\end{equation}
Normalisation is essential: raw stabilities are enormous and scale with point
count, so an absolute penalty never bites; $\Delta$ is a scale-free measure of how
much the node prefers to stay whole. Against it we apply a \textbf{soft size
pressure}, zero up to a single-sheet size and growing smoothly beyond. The node
stays whole when a logistic of \emph{($\Delta$ minus the size pressure)} exceeds
one half, else splits into its children; a kept node deselects everything below,
and the root is never a cluster. Splits can occur only at the tree's existing
branches---the weakest density seams---so the pressure decides only \textbf{whether}
to accept a seam split, never where to cut, and \textbf{recurses}, peeling a
three-way merge one seam at a time. A coherent sheet resists ($\Delta$ stays
high); a merge yields (its children are stable). The rule is deterministic, so the
labels carry no clustering randomness.

\paragraph{Recovering coverage without re-merging.} After selection the cores
occupy only part of the points; the peeled inter-sheet material is left as noise.
We fill it back \textbf{in the embedding only}: each unlabelled primitive takes
the majority identity of its nearest labelled cores---a Voronoi fill along the
boundaries the clustering drew, restoring coverage without dragging the cores
together. The distance travelled to reach those cores is kept as a signal (zero at
a core, large in deep filler): the embedding's contribution to per-voxel
confidence.

So heavily over-merged windows split into their true wraps while clean windows are
untouched---the correction acts only where the bias did damage. The trained
embedding and the sheets it recovers are shown in Figures~\ref{fig:embedding}
and~\ref{fig:sheets}.

\begin{figure}[h]
  \centering
  \figasset{sheets-with-embeddings.jpg}{\linewidth}{5cm}
  \caption{A 3-D PCA projection of the trained 8-D embedding for a real window,
  coloured by the probeom cluster assigned here: same-sheet primitives collapse
  into tight cores and the adjacent wraps lay out as an ordered ladder, with the
  propagated filler shaded by its distance-to-core (the low-confidence signal).
  The recovered sheets are shown alongside for reference.}
  \label{fig:embedding}
\end{figure}

\begin{figure}[h]
  \centering
  \figasset{sheets-with-sheets.jpg}{\linewidth}{5cm}
  \caption{The recovered sheets in a real window, each cluster drawn as its own
  solid-material medial surface in a distinct colour---the same clustering as the
  embedding figure above, now seen in the volume as separated papyrus sheets.}
  \label{fig:sheets}
\end{figure}

\section{From clusters to labels}
Per-primitive cluster identities become voxel labels by a geometric fit. For each
cluster we fit a \textbf{medial} surface plus a local thickness, rather than either
face, because a medial segments more cleanly and is less ambiguous. The fitted
sheets are rasterised to the two fields the refiner consumes:
\begin{itemize}[nosep]
  \item a \textbf{surface crest}, a soft three-way field marking surface (and where
  the label abstains); and
  \item a per-voxel \textbf{sheet identity} by nearest-medial-sheet (a Voronoi
  partition among the fitted sheets)---a genuine instance segmentation, not a binary
  mask, which is what lets us supervise the refiner's separation objective at all.
\end{itemize}

\section{Confidence and quality}\label{sec:confidence}
Nothing is thresholded or discarded at generation; every window carries honest
per-voxel and per-sheet uncertainty, so what to trust is decided later. Two
\textbf{regional} confidence fields are produced from independent evidence and
combined.

\paragraph{The intersection (geometric) confidence.} Where two rasterised medial
envelopes \emph{interpenetrate}, an \textbf{overlap} signal is written---a direct
symptom of a merge or mis-fit. We Gaussian-smooth it, \textbf{coverage-normalised}
(divided by the smoothed sheet mask, so air never dilutes it), and set
\begin{equation}
  \mathrm{conf}_{\text{intersection}} \;=\; 1 - \mathrm{smoothed\_overlap}
  \;\in[0,1],
\end{equation}
low where fitted sheets collide, high where they partition cleanly. It is a
separate channel and never modulates the crest.

\paragraph{The \code{sharp2} (embedding) confidence.} From the clustering's
\textbf{propagation distance}: each primitive carries $\mathrm{pd}$, the mean
embedding distance to the cores it was voted into (zero for a core, large for deep
filler or a residual merge). We normalise $\mathrm{pd}$ by its 90th percentile to
an uncertainty in $[0,1]$, \textbf{splat} it into the volume (summed uncertainty
and hit count per voxel), Gaussian-smooth both, and take the weighted average:
\begin{equation}
  \mathrm{conf}_{\text{sharp2}} \;=\; 1 - \langle\text{uncertainty}\rangle_{\text{smoothed}},
\end{equation}
low wherever ambiguous, far-propagated meshlets sit---the bridges, seams and
residual merges the geometry alone cannot see. (\code{sharp2} is a legacy channel
key.)

\paragraph{Combining them.} The channels merge as a \textbf{soft-OR of their
doubts}: with $u_i=1-\mathrm{conf}_i$, the combined confidence is
$1-\prod_i(1-u_i)$, algebraically the \textbf{product}
$\mathrm{conf}_{\text{intersection}}\cdot\mathrm{conf}_{\text{sharp2}}$. Either
channel's doubt pulls it down---a merge shows in the geometry, an ambiguous fill in
the embedding; the combined field on a hard window is shown in
Figure~\ref{fig:confidence}.

\paragraph{Per-sheet quality.} Separately, each fitted sheet is scored by the
agreement between its own surface normals and the fibre-field orientation. A
collapsed sheet---the fit forced through mush---shows poor congruence and scores
low. The score is \emph{stamped}, not acted on: the window is emitted regardless
and the collapse threshold ($\sim$0.80) is applied downstream. On our data it
flags roughly a third of windows as holding at least one collapsed sheet---an
honest measurement of where the method is not to be trusted.

\begin{figure}[h]
  \centering
  \figasset{sheets-with-confidence.jpg}{\linewidth}{5cm}
  \caption{A harder window than the clean case above---one dense with sheet
  intersections and wraps that do not fully separate in embedding space---with its
  recovered sheets coloured by a confidence
  heatmap (the combined field) instead of by cluster. The low-confidence
  intersections, seams, and residual merges read warm (red); confident sheet cores
  read cool, so the honest uncertainty falls exactly where the geometry is
  genuinely ambiguous.}
  \label{fig:confidence}
\end{figure}

\section{Augmentations}
Because the dominant error the refiner must undo is a \emph{merge}, each window
also emits controlled corruptions of itself, reusing primitives and clusters
otherwise discarded. A \textbf{merge} augmentation fuses adjacent sheets in the
embedding and paints the region low-confidence; a \textbf{split} augmentation
removes a sheet and re-partitions its space. These give paired examples of the
very failures---spurious welds and missing sheets---the refiner is trained to
repair, the affected regions marked low-confidence so the corruption is never
mistaken for signal.

\section{Validation: the signal is in the labels, not the loss}\label{sec:val1}
These pseudo-labels are only worth publishing if they carry sheet-\emph{separation}
signal a segmentation-derived model does not already have. The yardstick is
\textbf{m7}---the ScrollPrize \code{surface\_m7\_nnunet} detector, trained on
labels derived from existing segmentations. To measure what our labels add, we
fine-tuned m7 on our pseudo-labels alone (no human labels) and named it
\textbf{ablA} (ablation A).

A clean comparison must isolate the \emph{data} from the recipe, which forces a
decision about the loss. Because we cover whole windows sampled across the volume
rather than the neighbourhoods an existing segmentation delineates, our labels
carry far \textbf{more negative (non-sheet) voxels}, and a loss tuned to m7's
balance would be swamped. The ablation uses one suited to it, one job per term: a
\textbf{symmetric Focal--Tversky} term for precision, a \textbf{soft
skeleton-recall} term for coverage (recall measured only against the labelled
skeleton, so the abundant negatives cannot dilute it), and a \textbf{separation
penalty} suppressing predicted surface inside the inter-sheet gaps. This is
deliberately \emph{not} HercUNet's affinity/MALIS objective; the ablation holds
the loss \textbf{fixed} and varies only the labels, so any difference it exposes
is a property of the data.

The isolating experiment is \textbf{ablB}: fine-tune m7 with the \emph{same recipe
and loss} as ablA but on \textbf{m7's own labels}. If the loss did the work, ablB
would match ablA; if the signal lives in the data, ablA recovers structure ablB
cannot. It does (Figure~\ref{fig:ablAB}): ablB stays degenerate in the very regions
where ablA finds coherent sheet signal---the separation our labels encode is not
recoverable from segmentation-derived labels by a change of loss. On its own, ablA
also recovers coherent sheets deep in regions where the m7 baseline is degenerate
(Figure~\ref{fig:ablAsignal}).

\begin{figure}[tb]
  \centering
  \figasset{ablA-vs-ablB.jpg}{\linewidth}{5cm}
  \caption{ablA (trained on our labels) versus ablB (m7's own labels, our loss)
  on the same region. The sheet structure ablA recovers is absent from ablB---with
  the loss held fixed, the difference is carried by the labels, not the objective.}
  \label{fig:ablAB}
\end{figure}

\begin{figure*}[tb]
  \centering
  \figasset{ablA-signal.jpg}{\textwidth}{5cm}
  \caption{The public m7 volume on PHerc1447 (left) and the non-thresholded
  output of ablA (right). ablA recovers coherent sheets deep in a region where
  the segmentation-derived baseline is degenerate---evidence that the training
  labels, not the objective, supply the signal.}
  \label{fig:ablAsignal}
\end{figure*}

Two caveats. First, ablA tends to \textbf{over-merge} sheets in places---the
failure mode HercUNet (Part~\ref{part:refiner}) is built to address, shown later in
Figure~\ref{fig:merges}. Second, these are \textbf{pure pseudo-labels}: no cleanup,
no curation---the published corpus is exactly what the generator emits.

That the labels can be \emph{cleaned at all} follows from how they are built.
Because a window's sheets live in the 8-D embedding of \S\ref{sec:embedding}, a
correction is a \emph{cluster} operation, not voxel painting: an over-merged sheet
split by a cut, two fragments joined, a spurious cluster deleted, with the
medial-mesh fit and confidence re-derived from the corrected clustering. The
released tooling exposes exactly this---an interactive editor over the meshlet
embedding---so progressively cleaner label sets can be extracted from the same
corpus without re-running the pipeline.

\section{Reproducibility}
Generation is reproducible on two levels: every random draw (embedding init,
negative sampling, augmentation) is keyed to a seed derived from the scroll and
window coordinates, and the GPU reductions run in a deterministic mode, so
borderline labels do not flip at density seams from accumulation noise. Two
independent runs of a window reproduce every stored array exactly.

\textbf{The reproducible unit is the window, not the coordinate draw.} Each
window's master seed is
$\mathrm{md5}(\code{run\_id}\,|\,\text{scroll}\,|\,\text{level}\,|\,z,y,x)$, so it
is fully determined by \emph{which} window it is. The corpus was mined by many
parallel workers each drawing a disjoint seeded coordinate stream; reproduction is
per-window and content-addressed---regenerate the scroll + coordinates and you get
the same result, to within the equivalence of the code revisions the run spanned
(\emph{same method, same parameters}, not byte-identical). The full provenance is
catalogued separately, with the command that regenerates an equivalent corpus.

% =====================================================================
\part{Iterative Surface Refinement (HercUNet)}
\label{part:refiner}
% =====================================================================

We present the second stage: a learned refiner that turns the pseudo-labels of
Part~\ref{part:labels} into a dense, self-correcting predictor. It is a UNet based
on m7 generating a map $D(\mathrm{CT},\mathrm{prev})\to\mathrm{surface}$ applied
\emph{iteratively}---its own output at one pass becomes the \code{prev} input of
the next---so structure propagates one window at a time and incomplete predictions
are repaired against the raw CT, not a fixed label. Three ideas carry the design.
First, the network never has to remember where sheets are: a \textbf{carried
orientation field}, read from the structure tensor of the CT at every step,
supplies the geometric cue that separates touching wraps wherever the material
still holds two orientations---a cue any prediction that fused them would have
lost. Second, detection and \textbf{separation} are learned by different terms---a
symmetric detection substrate that finds surface, and a \textbf{constrained-MALIS
affinity objective} over per-sheet identity that grows a sheet along itself while
cutting it from its neighbours---so completing a sheet and merging two are no
longer the same gradient. Third, the iteration is \emph{taught}: a synthetic
\code{prev} channel, composed each epoch from label variants then cut and
fragmented, plus online DAgger conditioning on the model's \textbf{own} output,
closes the train/inference gap.

We are candid about the limit: completing and growing sheets works, but separating
two wraps \emph{already} merged---where the CT may carry no distinguishing
orientation---is not yet solved. We situate the work in learned tractography,
taking one idea---that a network can \emph{track} a structure through the raw
signal beyond the traces it was shown---to grow sheets into material the labels
never reached. That the pseudo-labels carry signal even where they appear to fail
is confirmed by the reverse experiment: \emph{removing} their noisy parts (dropping
the mushy-looking windows, substituting a segmentation-derived teacher) measurably
\textbf{degraded} the refiner. The uncurated, honestly-uncertain labels are the
asset.

\section{Setting: from pseudo-labels to a learned refiner}
The first stage produces, per window and unsupervised, two rasterised fields---a
thin surface \textbf{crest} and a per-voxel \textbf{sheet-instance identity}, each
with honest confidence---correct where the geometry is clean, uncertain where it
is not. But they are \emph{local} (decoded inside a 192-voxel window from its own
material) and only as complete as the primitives that produced them. They are not
a predictor: on a fresh volume they cannot be evaluated, do not fill the gaps a
window's tracing left, and do not improve with use.

The second stage distils them into a network that does all three. We take the
community's surface backbone (m7) and give it a second job---\textbf{refine} a
running surface estimate conditioned on the CT---trained so that applying it
repeatedly contracts toward a cleaner, more complete surface than any single pass
or window's pseudo-label contains. The refiner is an iteration-agnostic map
\begin{equation}
  D(\mathrm{CT},\,\mathrm{prev}) \to \mathrm{surface},
\end{equation}
where $\mathrm{prev}\in[0,1]$ is a surface-probability field: at inference the
model's own previous-pass softmax, at training a decoded pseudo-label crest
corrupted as in \S\ref{sec:teach}. Where $\mathrm{prev}\approx0$ (first pass, air,
a missed region) it predicts from CT alone; where \code{prev} is a crest it refines
it. The key invariant: \textbf{the target is always the clean base surface},
whatever \code{prev} shows. The model never learns a fixed
$\mathrm{prev}\to\mathrm{target}$ pairing---it learns to move \emph{any} state
toward the truth.

\section{Positioning: learning to track the sheets beyond the labels}
The first stage's primitives are literally \textbf{tractography}: each arm is a
streamline (an integral curve of a direction field, grown by fixed-step midpoint
integration), differing from fibre tracking only in \emph{which} field is
integrated (the sheet normal, not a diffusion ODF) and the constraints that keep
it on the sheet. We build on learned tractography not for labels---the first stage
supplies them---but because it has worked out how to make a network \emph{follow a
structure through the raw signal} past the traces it was trained on: exactly what
the refiner needs to grow a sheet into material the labels never reached.

Poulin et al.~\cite{poulin2019}, reviewing ML for tractography, characterise
learned tracking as \textbf{path-based} and \textbf{local-model-free}---stepping a
model along a trajectory rather than thresholding per-voxel---and note it can make
\textbf{non-local} decisions a local method cannot. That is what we want for
sheets. The supervised realisation, Poulin et al.~\cite{poulin2017}'s \emph{Learn
to Track} (recurrent nets predicting the next streamline direction), trains
against \textbf{reference streamlines} and so only reproduces the tracks it was
shown.

Th\'eberge et al.~\cite{theberge2021}'s \emph{Track-to-Learn} is the idea we
borrow. It casts tracking as reinforcement learning: an agent propagates a
streamline from the local signal, driven by a \textbf{reward} rather than
supervised targets, and so tracks from the signal itself, independently of any
reference tracks---the demonstration that a network can extend a structure through
the signal \textbf{further than its labels reach}.

We bring this to sheet tracing. The refiner is given the first stage's
pseudo-labels---it is not label-free---but is deliberately \textbf{not bounded by
them}. The affinity and material-growth terms of \S\ref{sec:loss} and the
iteration of \S\ref{sec:teach} teach it to \emph{track the sheet through the CT
itself}: to complete a sheet across a gap and extend it into unlabelled material,
exactly as a signal-driven tracker carries a streamline beyond its reference
tracks. The labels say what a sheet \emph{is}; the CT says how far it \emph{goes}.

\begin{figure*}[tb]
  \centering
  \figasset{grow-over-m7_PHerc0175A_z5568_y4224_x2304.jpg}{\textwidth}{4.5cm}
  \caption{We fed the output of m7 to our model, which completes the gaps
  coherently using the idea of learned tractography. From left to right: the CT
  scan, the output of m7, one iteration of our model taking m7's output as its
  \code{prev} input, and the difference between \code{prev} and our model's
  output.}
  \label{fig:growm7}
\end{figure*}

\section{The refiner as an iterated map}
At inference the refiner runs as a small number of \textbf{double-buffered Jacobi
passes}. Pass~0 gets $\mathrm{prev}=0$ everywhere and predicts from CT alone; each
later pass reads the whole previous pass's surface-probability volume as its
\code{prev} and writes a fresh one. Information propagates roughly one window per
pass---a deliberately \emph{local} refinement, not a global solver---which is all
the completion the problem needs: a sheet a cold pass detected only in patches is,
pass by pass, joined along itself.

Between passes the grid shifts by half a stride (\S\ref{sec:infer}) so the seams
of one pass fall at the confident centres of the next and are healed---the
\emph{spatial} counterpart of the self-conditioning taught at training
(\S\ref{sec:dagger}), where the model meets a \code{prev} whose octants came from
different neighbouring windows and must reconcile them. The network is a stock
nnU-Net ResEnc backbone, warm-started from m7 (\S\ref{sec:affinity}), so a model
good at \emph{detecting} surface is re-tasked to \emph{refine} it.

\section{The carried orientation field}\label{sec:orient}
A refiner conditioning on its own output risks a specific failure: once two
touching wraps are fused in \code{prev}, nothing in a purely intensity-driven
input tells the next pass they were ever two. The fix is to give the network, at
every pass, the one cue that survives the fusion---the \textbf{local orientation
of the material}---read from the CT, not from any prediction.

We append to the $[\mathrm{CT},\mathrm{prev}]$ input \textbf{six orientation
channels}: the entries of the unit-trace \textbf{structure tensor} of the CT,
$[J_{zz},J_{yy},J_{xx},J_{zy},J_{zx},J_{yx}]$, from the same
derivative-of-Gaussian-and-integrate operator as the first stage, normalised by
trace so only the \emph{shape} of the orientation---not intensity---is carried. It
is \textbf{recomputed inside the forward pass from the augmented CT}, never stored,
so whatever augmentation the CT underwent the orientation is automatically
consistent, and it is \textbf{sign-free} (the tensor is even under
$\nhat\to-\nhat$). We read it from the CT because, \emph{where the material retains
two orientations}, the CT keeps them even when a detector fused the wraps into one
crest, whereas an orientation read off the prediction would inherit that fusion;
the separating cue is thus still in the input though the prediction lost it.
(Reading from the CT also lets the field transform correctly under reflections---the
basis for the test-time mirror augmentation of \S\ref{sec:infer}.)

This is where the method is weakest, facing \textbf{two failure modes}, neither
solved. First, the cue is \textbf{not always there}: in the most compressed
regions two wraps press close enough that the material has no locally
distinguishable orientation---the structure tensor sees one coherent direction,
as for a genuine sheet---and the input carries no separating signal; only
\textbf{non-local} evidence (an adjacent window where the two \emph{do} separate,
carried across passes) could recover it, and doing that reliably is still ahead of
us. Second, even where the cue \emph{is} present, getting the model to \textbf{act}
on it and pull an already-merged pair apart is the hardest thing we ask, and the
current model does not do it reliably. The constrained-MALIS objective of
\S\ref{sec:malis} and this orientation input target both. Growing and completing
sheets through gaps works; \textbf{separating already-merged sheets is an open
problem we are actively working on}, not a capability we report.

\section{The affinity head and instance separation}\label{sec:affinity}
The backbone is \textbf{m7}, a single-channel $[\mathrm{CT}]\to\mathrm{surface}$
model and a strong \emph{detector}---which is why we build on it. Warm-starting
copies its encoder and decoder verbatim and \textbf{stem-expands} its first
convolution from one input channel to eight (CT weight on channel~0; \code{prev}
and the six orientation channels start at zero), so the refiner begins predicting
exactly what m7 predicted and departs from there.

To the shared decoder we attach a second, full-resolution \textbf{affinity
head}---a $1\times1\times1$ convolution off the last decoder feature map---so the
network emits two things at once: the surface segmentation \code{seg}, and a
field of per-voxel \textbf{affinities} $\mathrm{aff}\in[0,1]$, one channel per
spatial offset, each predicting whether a voxel and its offset neighbour belong
to the \emph{same sheet}. Detection and instance-separation read from the same
features but are supervised by different losses (\S\ref{sec:loss}), so surface
presence and sheet membership are decoupled.

\section{The loss stack}\label{sec:loss}
Three families act together: a detection substrate on \code{seg}, a topological
separation objective on \code{aff}, and a data-driven growth term on \code{aff}.
Only the first is inherited from the single-pass ablation; the other two make the
refiner an \emph{instance} model, not a thicker detector.

\subsection{Detection substrate: symmetric Focal--Tversky, separation, skeleton-recall}
The detection loss is the same stack as the single-pass model, so the two differ
only in machinery, not in the definition of surface:
\begin{equation}
  \resizebox{\columnwidth}{!}{$\displaystyle
  \mathrm{DeepSup}\big(\mathrm{SkelRecall}\big(\mathrm{Separation}\big(\mathrm{DC\_CE}(\mathrm{CE}{+}\text{Focal--Tversky})\big)\big)\big)$},
\end{equation}
evaluated with \code{ignore} masked out. The core is nnU-Net's \code{DC\_and\_CE}
with Dice replaced by a \textbf{Focal--Tversky} term,
$\mathrm{TI}=\mathrm{TP}/(\mathrm{TP}+\alpha\,\mathrm{FP}+\beta\,\mathrm{FN})$,
$\mathrm{loss}=(1-\mathrm{TI})^\gamma$. Crucially it is \textbf{symmetric}
($\alpha=\beta=0.5$): a recall-tilted Tversky and the skeleton-recall term below
would both push activations up and collapse the prediction into mush, so we make
Tversky purely the \textbf{precision} term and leave recall and anti-merge to the
dedicated terms. Recall comes from a \textbf{skeleton-recall} term (clDice~\cite{cldice}
soft-skeleton recall,
$\sum(P_{\text{surf}}\cdot\mathrm{skel}(\mathrm{GT}))/\sum\mathrm{skel}(\mathrm{GT})$,
ignore-masked), rewarding coverage of the thin GT medial without penalising
anything off it. Anti-merge comes from a \textbf{separation penalty}: the
morphological closing of the surface minus the surface, restricted to background
voxels, identifies the thin inter-sheet \textbf{gaps} and penalises predicted
surface inside them. Because \code{ignore} and \code{surface} are excluded by
construction, it never fires inside a region the label already marked uncertain.

\subsection{Constrained MALIS on the affinity field}\label{sec:malis}
The separation penalty is local. Long-range membership---that two crests a
millimetre apart belong to the same wrap---is learned by a \textbf{constrained
MALIS} objective (Turaga et al.~\cite{turaga2009}) in its two-pass form (Funke et
al.~\cite{funke2018}), over the first stage's per-sheet identity.

MALIS scores an affinity field by its effect on \textbf{segmentation topology}:
the connectedness of two voxels is the \emph{maximin} edge---the weakest affinity
on the strongest path between them. A \textbf{positive} pass considers within-sheet
pairs and pushes their maximin edge up (growing and completing each sheet,
long-range by construction); a \textbf{negative} pass fuses each GT sheet, finds
the maximin edge \emph{between} sheets, and pushes it down (cutting the one weak
link that would merge them). Maximin edges are found exactly with Kruskal over a
union--find, so no external MALIS dependency is needed; the GT is the mesh-Voronoi
instance id, background and \code{ignore} excluded. Affinities are predicted at
short, mid and long offsets (nearest neighbours plus 3 and 9 along each axis),
evaluated on a random sub-crop per step to bound cost. The weight is
\textbf{ramped} from zero, so the network learns to \emph{detect} on the
warm-started substrate before it learns to \emph{separate}.

\subsection{CT-material growth and air-suppression}
MALIS is supervised only where the sparse instance labels reach. To complete
sheets through the large \textbf{unlabelled} stretches between windows---and forbid
growth into air---a third term drives the affinity field from the CT with no
labels. A material scaffold $m(x)=\sigma((\mathrm{CT}-\tau)/s)$ reads high on
papyrus, low in air. An edge is asked to \textbf{grow} ($\mathrm{affinity}\to1$)
with weight $m(u)\,m(v)\cdot\mathrm{along}$,
$\mathrm{along}=1-\hat{\delta}^{\mathsf T}J\hat{\delta}$ high when the offset runs
\emph{along} the tangent, so growth completes sheets \textbf{in-plane through
material}, while across-normal growth is left to MALIS so stacked sheets are never
bridged. An edge with either endpoint in air is asked toward \textbf{zero}
($1-m(u)m(v)$), a hard no-grow through voids. Both weights are data-derived and
detached; the term is ramped in. This lets the refiner extend a sheet across a gap
the pseudo-label never covered without hallucinating in air---checked per-epoch
with a validation cross-section overlaying prediction, label and material gate on
the CT.

\section{Teaching iteration: the synthetic \code{prev} channel}\label{sec:teach}
The refiner is only useful if applying it repeatedly improves the result---a
property of \emph{training}, not architecture. The danger is a train/inference
mismatch: a model that only sees clean, label-derived \code{prev} never learns to
repair the incomplete, slightly-wrong \code{prev} it feeds itself at inference.
Two mechanisms close the gap, both under one discipline: whatever the \code{prev},
the \textbf{target is the clean base surface}, so every corruption below is
something the loss asks the model to fix.

\subsection{Octant composition}
Each window stores, besides its clean base crest, a set of \textbf{candidate}
fields---merge and split \emph{augmentation variants}, each a valid decoded crest.
Every epoch a transform composes \code{prev} afresh from these. With some
probability the whole window is zeroed (the bootstrap / first-pass regime, teaching
cold prediction from CT); otherwise it is split into $2\times2\times2$
\textbf{octants} at a jittered centre, each filled from an \textbf{independently
sampled} candidate (zero, clean base, or an augmentation). The result is a
plausible-but-inconsistent \code{prev} disagreeing across a plus-shaped seam at the
centre---exactly the inference geometry, where a half-shifted window's octants were
written by different previous-pass windows. The model is thereby taught to trust
the CT over an inconsistent prior and reconcile the seam. Crucially the composition
only \emph{selects and stitches valid candidate fields}; it never invents geometry.

\subsection{Fragmentation: masked-autoencoding on \code{prev}}
A refiner trained only to denoise a mostly-complete \code{prev} will tidy its input
but not \textbf{re-connect} what a previous pass left broken. To teach
gap-bridging we apply structured dropout to \code{prev} \textbf{only}, leaving the
target and loss mask untouched, so the clean label still supervises every cut voxel
as surface and the loss penalises failing to re-fill it---that penalty \emph{is}
the bridge-the-gap signal. Two modes alternate: \textbf{oriented slabs and boxes}
sever a sheet into disconnected sections (within-window completion, both ends
visible---interpolation); and a \textbf{half-space cut} removes all \code{prev} on
one side of a random plane, leaving a sheet entering from a single edge and forcing
\textbf{extrapolation} into unknown territory---the cross-window skill the
half-shifted tiling demands. This is masked-autoencoding on the prior channel, and
it converts an iteration that merely cleans into one that also grows.

\subsection{Online DAgger: conditioning on the model's own output}\label{sec:dagger}
Octant composition and fragmentation still start from \emph{label-derived} crests.
The remaining mismatch---that at inference \code{prev} is the model's \textbf{own},
differently-flawed output---is closed by online DAgger~\cite{dagger}. With some
probability a step first runs a few \code{no\_grad} warm-up forwards, each feeding
the model's own output back in as \code{prev} (the inference recurrence), then takes
one graded forward on that self-generated \code{prev}. The warm-ups hold no
activations, so the step costs little more than a single-pass run, and more of them
teach repair at deeper iteration depths cheaply; there is no back-propagation
through the recurrence. Because the warm-up's start still comes from the
octant/fragment composition, the self-outputs to repair are diverse---cold,
label-like, and cut---covering the full range of states a real iteration passes
through.

We keep split-teaching honest with one asymmetry: \textbf{split} variants, which
pair a split \code{prev} with a single-sheet target and would teach \emph{merging},
are dropped by default; base and \textbf{merge} variants, which teach separation,
are kept. Our failure mode is over-merging, so the augmentation distribution is
chosen not to reinforce it.

\section{Blending m7 in the compressed regions}\label{sec:blend}
Our labels are strongest where the material is well-formed; they weaken in the most
\textbf{compressed} regions, where wraps press to contact and even the orientation
cue is marginal, and there m7 is often the better teacher. We therefore
\textbf{mine} a modest set of m7 pseudo-labels from compressed regions and blend
them in, train-only, at a small fraction of the corpus.

Because the refiner is warm-started from m7 (\S\ref{sec:affinity}), this blend is
best read as a \textbf{rehearsal} signal against catastrophic forgetting, not new
supervision. In the compressed regions the network could segment \emph{as m7} the
day training began; a steady diet of m7 labels there reminds it---\emph{you were
able to do this; do not forget how}---while every other window teaches our thinner,
instance-aware crest. It does not hand the network a new skill so much as stop it
\emph{unlearning} an old one exactly where our own labels fall silent.

The recipe is conservative and is the source of truth shared with the cleanup of
\S\ref{sec:signal}. From an m7 surface binary we compute m7's \textbf{own
normal-coherence} (the structure tensor of the smoothed mask) and keep as positive
\code{surface} only the coherence-gated crest ($\text{m7}\wedge\mathrm{coh}>\tau$);
the incoherent remainder is marked \code{ignore}, the rest background. The kept
label is m7's coherent \textbf{face} at native thickness---not a medial---a
different geometric object from our thin \dgrad\ crest, and that difference is
itself the labels-vs-loss variable the first stage isolates. These m7 cases carry
\textbf{no instance identity} (MALIS is silent on them) and no candidate crests
(their \code{prev} bootstraps to zero, training \textbf{cold} detection), and are
tagged by source so training logs the two populations (\code{our} / \code{m7})
separately. The blend is small by design---to cover the compressed regions our
labels cannot, not to let a segmentation-derived teacher dominate.

\begin{figure}[tb]
  \centering
  \figasset{m7-vs-hercunet-merges.jpg}{\linewidth}{5cm}
  \caption{The mined m7 labels help \emph{a lot} with retaining separation in
  semi-compressed regions. In a half-compressed region---where our own
  pseudo-labels are thin and a refiner trained without the blend tends to
  over-merge---HercUNet keeps the adjacent wraps apart (shown against m7). The
  small, train-only m7 blend supplies the separation supervision there, exactly
  where our own labels reach only thinly.}
  \label{fig:merges}
\end{figure}

\section{The herculabels carry signal even where they appear to fail}\label{sec:signal}
Because m7 is the stronger teacher in the most compressed regions
(\S\ref{sec:blend}), the tempting next step is to let it correct us wholesale:
discard the windows where our label looks mushy and, where m7 clearly beats us,
\textbf{replace} our label with m7's coherent one. We built exactly that---a
cleanup that, only where m7's advantage exceeded a margin, either \textbf{dropped}
the case or \textbf{replaced} our surface with m7's coherence-gated label;
competitive-or-better windows were left untouched.

The result is the important negative finding of this work: the cleaned corpus
trained a refiner \textbf{no better than the uncleaned one}, regressing to the
ablation that drops uncertain-region supervision. The mushy-looking labels we
removed were carrying signal. Two losses compounded: dropping the mushy windows removed
exactly the hard, compressed, ambiguous material the refiner most needed, leaving
the easy windows it could already do; and replacing our failures with m7 overwrote
our rough-but-CT-grounded reading with a segmentation-derived answer that, by our
thesis, cannot recover what the CT still carries. Only a small fraction was
touched, and even that erased the gain (Figure~\ref{fig:cleanup}).

The lesson does \textbf{not} contradict \S\ref{sec:blend}. m7 is valuable as
\emph{additional} coverage where our labels are absent or thin, not as a
\emph{substitute} where we already have them, however rough: the herculabels carry
real signal in exactly the noisy places they appear to fail, and marking those
regions uncertain and keeping them beats deleting or overwriting them. We
therefore train on the \textbf{uncurated} corpus, keep the m7 blend as purely
additive coverage, and never let cleanup drop or replace our own hard cases.

\begin{figure*}[tb]
  \centering
  \figasset{cleaned-labels-signal-loss.jpg}{\textwidth}{5cm}
  \caption{The cleaned labels were also carrying signal. A 2\,mm XY cross-section
  of PHerc1447: \textbf{m7} (blue), and two refiners trained with the \emph{same}
  recipe but different label sets --- \textbf{ours on the cleaned/curated corpus}
  (yellow; mushy-looking windows dropped, failures m7-substituted) and
  \textbf{HercUNet} on the full uncurated corpus (magenta). The cleaned-corpus
  model loses the coherent sheet signal HercUNet recovers here: removing the
  mushy-looking labels removed real signal (\S\ref{sec:signal}).}
  \label{fig:cleanup}
\end{figure*}

\section{Inference: iterative passes and the seam-free blend}\label{sec:infer}
Inference runs the refiner for a small number of passes over a region, streaming
the CT through the shared data layer and writing each pass's surface probability
as an OME-Zarr so the iterations are inspectable. \textbf{Within each pass},
overlapping windows are merged by a Gaussian importance map---the same seam-free
stitching the single-pass detector uses---so no within-pass grid artefact
survives. \textbf{Between passes} the grid shifts by half the \emph{stride}, so
each pass's window centres land on the previous pass's overlap-centres: the seams
of one pass are re-predicted at the confident centres of the next and healed. This
is the \textbf{iterate $\to$ blend $\to$ iterate} loop, at a compute cost set by
the overlap. Each output tile is chunk-aligned and written exactly once, so a run
is reproducible and splits across workers with no tile written twice.

\section{Evaluation: why we report no single number}\label{sec:eval}
We deliberately publish no Dice, VOI, or leaderboard figure. This is not an
omission we intend to quietly patch---the reasons are structural, and stating them
is more honest than a number we know would mislead.
\begin{itemize}[nosep]
  \item \textbf{Leakage.} m7 is trained on the full published label set (700+
  Kaggle surface labels). Validating on those, or on anything mined from m7,
  scores a model against data its teacher has already seen; there is no held-out
  ground truth in our domain independent of m7.
  \item \textbf{Dice measures label thickness, not detection.} Dice is dominated
  by the \emph{thickness} of the reference---one voxel thicker moves the score
  substantially, whether or not the sheet was found. Against labels of
  uncontrolled thickness it reports label metrology, not surface correctness.
  \item \textbf{Faces vs.\ medials---a category difference, not an error.} Every
  label available to us annotates the \textbf{face} at native thickness; our
  pseudo-labels annotate the \textbf{medial} crest, a few voxels inboard
  (Figure~\ref{fig:facemed}). An overlap metric charges that constant offset as a
  miss, when it is the correct centre of the same sheet.
  \item \textbf{Human smoothness bias.} Our labels track the material where it
  genuinely bends, pinches, and delaminates; human labels are known to
  smooth---annotators draw flatter sheets than the papyrus is. Scoring against them
  would penalise our labels \emph{for tracking the material more faithfully than a
  human}.
\end{itemize}

\paragraph{The case for the medial.} We defend it directly: \emph{it is better to
have a sheet detection whose face lives flat at a few voxels' offset than a
detection that is flat at zero offset only in some regions and waves harshly
everywhere.} A medial a consistent few voxels inboard is trivially recoverable
(one dilation to either face) and, more importantly, \emph{coherent}: it does not
wander. The alternative---a face-fit that is exact where the material is easy and
buckles where it is hard---is what actually breaks downstream flattening and
reading, because its error is \emph{unpredictable} rather than a constant,
correctable offset.

\paragraph{What we show instead.} In place of a scalar we show the prediction
itself---across the iterative passes and on volumes m7's labels never covered
(Figures~\ref{fig:banner}, \ref{fig:growm7} and~\ref{fig:0800})---and invite the
reader to judge what matters: coherence, completion of broken sheets, and
separation of touching wraps, properties a single overlap number against biased,
differently-thick, leakage-prone labels cannot report.

\paragraph{A note on generalisation---a hypothesis, not a result.} The refiner is
trained on \textbf{4{,}000+ HercuLabels windows across 22 carbonised scrolls}, not
the labelled neighbourhoods of a single corpus. On the standard argument that
training breadth improves out-of-distribution behaviour, one would expect better
cross-volume generalisation than a per-corpus detector, and the cross-scroll
results are at least consistent with that. But with no independent held-out ground
truth we cannot quantify it, so we offer it as a reasoned expectation, not a
measured claim.

\begin{figure*}[tb]
  \centering
  \figasset{PHerc0800-m7-vs-ours.jpg}{\textwidth}{5cm}
  \caption{m7 (blue) versus our refiner (magenta) on a 1\,cm cross-section of
  PHerc0800.}
  \label{fig:0800}
\end{figure*}

% =====================================================================
\begin{thebibliography}{9}

\bibitem{poulin2019}
P.~Poulin, D.~J\"orgens, P.-M.~Jodoin, and M.~Descoteaux.
\newblock Tractography and machine learning: Current state and open challenges.
\newblock \emph{Magnetic Resonance Imaging}, 64:37--48, 2019.

\bibitem{poulin2017}
P.~Poulin, M.-A.~C\^ot\'e, J.-C.~Houde, L.~Petit, P.~F.~Neher, K.~H.~Maier-Hein,
  H.~Larochelle, and M.~Descoteaux.
\newblock Learn to Track: Deep Learning for Tractography.
\newblock In \emph{Medical Image Computing and Computer-Assisted Intervention
  (MICCAI)}, 2017.

\bibitem{theberge2021}
A.~Th\'eberge, C.~Desrosiers, M.~Descoteaux, and P.-M.~Jodoin.
\newblock Track-to-Learn: A general framework for tractography with deep
  reinforcement learning.
\newblock \emph{Medical Image Analysis}, 2021.

\bibitem{turaga2009}
S.~C.~Turaga, K.~L.~Briggman, M.~Helmstaedter, W.~Denk, and H.~S.~Seung.
\newblock Maximin affinity learning of image segmentation.
\newblock In \emph{Advances in Neural Information Processing Systems (NeurIPS)},
  2009.

\bibitem{funke2018}
J.~Funke, F.~D.~Tschopp, W.~Grisaitis, A.~Sheridan, C.~Singh, S.~Saalfeld, and
  S.~C.~Turaga.
\newblock Large scale image segmentation with structured loss based deep learning
  for connectome reconstruction.
\newblock \emph{IEEE Transactions on Pattern Analysis and Machine Intelligence},
  41(7):1669--1680, 2019.

\bibitem{weickert}
J.~Weickert.
\newblock \emph{Anisotropic Diffusion in Image Processing}.
\newblock Teubner, Stuttgart, 1998.

\bibitem{vicreg}
A.~Bardes, J.~Ponce, and Y.~LeCun.
\newblock VICReg: Variance-Invariance-Covariance Regularization for
  Self-Supervised Learning.
\newblock In \emph{International Conference on Learning Representations (ICLR)},
  2022.

\bibitem{hdbscan}
R.~J.~G.~B.~Campello, D.~Moulavi, and J.~Sander.
\newblock Density-based clustering based on hierarchical density estimates.
\newblock In \emph{Pacific-Asia Conference on Knowledge Discovery and Data Mining
  (PAKDD)}, 2013.

\bibitem{cldice}
S.~Shit, J.~C.~Paetzold, A.~Sekuboyina, et al.
\newblock clDice --- A Novel Topology-Preserving Loss Function for Tubular
  Structure Segmentation.
\newblock In \emph{IEEE/CVF Conference on Computer Vision and Pattern Recognition
  (CVPR)}, 2021.

\bibitem{dagger}
S.~Ross, G.~J.~Gordon, and J.~A.~Bagnell.
\newblock A reduction of imitation learning and structured prediction to no-regret
  online learning.
\newblock In \emph{International Conference on Artificial Intelligence and
  Statistics (AISTATS)}, 2011.

\end{thebibliography}

\end{document}
